Source : calculs de l'auteur (fiscalite/is/calculs\_note\_is.py, mécanique de fiscalite/is/simulation\_is.py).\par Note : régime permanent, avant la baisse finale du taux. La déduction immédiate ne change pas le total déduit : sa valeur est celle de l'avance d'impôt, actualisée à 5 \% (taux des entreprises), pour une année d'achats ; elle vaut chaque année 3,8 \% des achats d'équipements. La première année, l'avance en trésorerie est bien plus grande : 34 800 \euro{} pour la PME industrielle, 2 300 \euro{} pour la PME de services, 35 M\euro{} pour le grand groupe. Les PME achètent sous le plafond de 1 M\euro{} et déduisent 100 \% dès 2027. Grand groupe : chiffre d'affaires d'au moins 3 Md\euro{}, sans CIR fractionné ni boîte à brevets ; impôt actuel 129 M\euro{} (taux de 25,8 \%), 181 M\euro{} en 2026 avec la contribution exceptionnelle.
